% Part III — The Trust Algebra (Slides 41–60)
\providecommand{\jugmark}[1]{\textbf{\color{jugeoBlue}#1}}

% ---------------------------------------------------------------------------
% Slide 41: Section divider
% ---------------------------------------------------------------------------
\begin{frame}
\frametitle{}
\vfill
\begin{center}
  {\Huge\color{jugeoBlue}\textbf{Part III}}\\[12pt]
  {\LARGE\color{jugeoDark}Trust --- The Key Innovation}\\[20pt]
  \begin{tikzpicture}
    \draw[jugeoBlue, line width=2pt] (-3,0) -- (3,0);
  \end{tikzpicture}\\[12pt]
  {\large\color{gray}\itshape Not all evidence is created equal}
\end{center}
\vfill
\end{frame}

% ---------------------------------------------------------------------------
% Slide 42: The Problem with Binary Trust
% ---------------------------------------------------------------------------
\begin{frame}
\frametitle{The Problem with Binary Trust}
\begin{center}
\begin{tikzpicture}[
  box/.style={draw, thick, rounded corners, minimum width=5cm,
              minimum height=1cm, font=\large, align=center}
]
  \node[box, fill=jugeoGreen!20] (yes) {\textcolor{jugeoGreen}{\textbf{\checkmark}}~Type-checks $\Rightarrow$ \textbf{100\% trusted}};
  \node[box, fill=jugeoRed!15, below=0.8cm of yes]
        (no)  {\textcolor{jugeoRed}{\textbf{\texttimes}}~Doesn't $\Rightarrow$ \textbf{0\% trusted}};
  \node[draw=jugeoOrange, dashed, thick, rounded corners,
        minimum width=5cm, minimum height=1cm,
        below=0.8cm of no, font=\large, align=center]
        (mid) {\textcolor{jugeoOrange}{???}~\emph{Nothing in between}};
\end{tikzpicture}
\end{center}
\vfill
\begin{itemize}
  \item LEAN / Coq / F*: proof accepted or rejected --- that's it
  \item Real-world evidence has \jugmark{degrees of confidence}
  \item Z3 proved 80\%, tests covered 15\%, expert reviewed 5\% --- where does that go?
\end{itemize}
\end{frame}

% ---------------------------------------------------------------------------
% Slide 43: Analogy — Courtroom Evidence
% ---------------------------------------------------------------------------
\begin{frame}
\frametitle{Analogy --- Courtroom Evidence}
\begin{center}
\begin{tikzpicture}[
  ev/.style={draw, thick, rounded corners, minimum width=3.4cm,
             minimum height=0.8cm, font=\small, align=center, text=white},
  arr/.style={-{Stealth[length=2.5mm]}, thick, gray}
]
  \node[ev, fill=jugeoGreen!80]                     (dna)  {DNA / Forensics};
  \node[ev, fill=jugeoBlue,  below=0.4cm of dna]    (eye)  {Eyewitness Testimony};
  \node[ev, fill=jugeoOrange, below=0.4cm of eye]   (circ) {Circumstantial Evidence};

  \node[right=0.6cm of dna,  font=\small, text=jugeoGreen!80]  {High trust};
  \node[right=0.6cm of eye,  font=\small, text=jugeoBlue]      {Medium trust};
  \node[right=0.6cm of circ, font=\small, text=jugeoOrange]    {Lower trust};

  \draw[arr] (dna)  -- (eye);
  \draw[arr] (eye)  -- (circ);
\end{tikzpicture}
\end{center}
\vfill
\begin{itemize}
  \item A jury doesn't rely on a single witness
  \item They \textbf{weigh} multiple channels of evidence
  \item The verdict reflects the \emph{combined} picture
  \item \jugmark{JuGeo works the same way for code verification}
\end{itemize}
\end{frame}

% ---------------------------------------------------------------------------
% Slide 44: JuGeo's Evidence Channels
% ---------------------------------------------------------------------------
\begin{frame}
\frametitle{JuGeo's Evidence Channels}
\begin{center}
\begin{tikzpicture}[
  ch/.style={draw, thick, rounded corners, minimum width=3.6cm,
             minimum height=0.7cm, font=\small, align=center},
  arr/.style={-{Stealth[length=2mm]}, thick, jugeoBlue!50}
]
  \node[ch, fill=jugeoGreen!15]                       (k) {\textbf{1.}~Formal Proof (kernel)};
  \node[ch, fill=jugeoGreen!10, below=0.25cm of k]   (z) {\textbf{2.}~SMT Solver (Z3)};
  \node[ch, fill=jugeoBlue!10,  below=0.25cm of z]   (t) {\textbf{3.}~Runtime Tests};
  \node[ch, fill=jugeoOrange!10,below=0.25cm of t]   (a) {\textbf{4.}~AI / Copilot Suggestions};
  \node[ch, fill=jugeoOrange!8, below=0.25cm of a]   (h) {\textbf{5.}~Human Oracle Review};
  \node[ch, fill=jugeoLight,    below=0.25cm of h]   (c) {\textbf{6.}~Composed Evidence};

  \node[draw=jugeoDark, very thick, rounded corners,
        fit=(k)(z)(t)(a)(h)(c), inner xsep=12pt, inner ysep=6pt,
        label={[font=\small\bfseries, text=jugeoDark]above:Evidence Sources}] {};
\end{tikzpicture}
\end{center}
\vfill
\begin{itemize}
  \item Each channel produces evidence with a \textbf{different trust level}
  \item JuGeo tracks \emph{which} channel justified \emph{which} claim
  \item \jugmark{No channel is ignored; no channel is unconditionally trusted}
\end{itemize}
\end{frame}

% ---------------------------------------------------------------------------
% Slide 45: The Seven Trust Levels
% ---------------------------------------------------------------------------
\begin{frame}
\frametitle{The Seven Trust Levels}
\begin{center}
\begin{tikzpicture}[
  lvl/.style={draw, thick, minimum width=6.5cm, minimum height=0.55cm,
              font=\small\ttfamily, align=center}
]
  \node[lvl, fill=jugeoGreen!25]                       (v) {VERIFIED\_PROOF};
  \node[lvl, fill=jugeoGreen!15, below=0pt of v]       (s) {SOLVER\_DISCHARGED};
  \node[lvl, fill=jugeoBlue!15,  below=0pt of s]       (r) {RUNTIME\_WITNESSED};
  \node[lvl, fill=jugeoOrange!15,below=0pt of r]       (o) {ORACLE\_PROPOSED};
  \node[lvl, fill=jugeoOrange!10,below=0pt of o]       (c) {COPILOT\_SUGGESTED};
  \node[lvl, fill=gray!10,       below=0pt of c]       (u) {UNVERIFIED};
  \node[lvl, fill=jugeoRed!15,   below=0pt of u]       (x) {CONTRADICTED};

  % arrow on the left
  \draw[{Stealth[length=3mm]}-{Stealth[length=3mm]}, very thick, jugeoBlue]
    ([xshift=-1.2cm]x.south west) -- ([xshift=-1.2cm]v.north west);
  \node[rotate=90, font=\small\bfseries, text=jugeoBlue]
    at ([xshift=-2cm]r.west) {Trust};

  % labels on the right
  \node[right=0.3cm of v, font=\scriptsize, text=jugeoGreen] {gold standard};
  \node[right=0.3cm of x, font=\scriptsize, text=jugeoRed]   {evidence \emph{against}};
\end{tikzpicture}
\end{center}
\vfill
\begin{center}
\small Every claim in JuGeo carries exactly one of these levels.\\
\jugmark{Higher is stronger; lower is weaker.}
\end{center}
\end{frame}

% ---------------------------------------------------------------------------
% Slide 46: Trust Level — VERIFIED_PROOF
% ---------------------------------------------------------------------------
\begin{frame}
\frametitle{Trust Level --- \texttt{VERIFIED\_PROOF}}
\begin{exampleblock}{The Gold Standard}
A \textbf{complete formal proof} checked by JuGeo's kernel.\\
Like a mathematical theorem --- \textbf{certain} (relative to the kernel's axioms).
\end{exampleblock}
\vspace{0.3cm}
\begin{center}
\begin{tikzpicture}[
  box/.style={draw, thick, rounded corners, minimum width=2.8cm,
              minimum height=0.8cm, font=\small, align=center}
]
  \node[box, fill=jugeoLight]                       (claim) {Claim};
  \node[box, fill=jugeoGreen!20, right=1.5cm of claim] (kern) {Kernel checks\\every step};
  \node[box, fill=jugeoGreen!30, right=1.5cm of kern]  (ok)   {\textcolor{jugeoGreen}{\checkmark}~\texttt{VERIFIED}};

  \draw[-{Stealth}, thick, jugeoBlue] (claim) -- (kern);
  \draw[-{Stealth}, thick, jugeoGreen] (kern) -- (ok);
\end{tikzpicture}
\end{center}
\vfill
\begin{itemize}
  \item No external dependency --- pure logical deduction
  \item Analogous to the verdict of a flawless DNA match
  \item \jugmark{Highest trust level in the system}
\end{itemize}
\end{frame}

% ---------------------------------------------------------------------------
% Slide 47: Trust Level — SOLVER_DISCHARGED
% ---------------------------------------------------------------------------
\begin{frame}
\frametitle{Trust Level --- \texttt{SOLVER\_DISCHARGED}}
\begin{block}{SMT Solver Proved It}
Z3 (or another SMT solver) automatically proved the claim.\\
\textbf{Very high confidence} --- but depends on Z3's own correctness.
\end{block}
\vspace{0.3cm}
\begin{center}
\begin{tikzpicture}[
  box/.style={draw, thick, rounded corners, minimum width=2.6cm,
              minimum height=0.8cm, font=\small, align=center}
]
  \node[box, fill=jugeoLight]                        (claim) {Claim};
  \node[box, fill=jugeoBlue!20, right=1.5cm of claim]  (z3)   {Z3 searches\\for proof};
  \node[box, fill=jugeoGreen!20, right=1.5cm of z3]    (ok)   {\textcolor{jugeoGreen}{\checkmark}~\texttt{DISCHARGED}};

  \draw[-{Stealth}, thick, jugeoBlue] (claim) -- (z3);
  \draw[-{Stealth}, thick, jugeoGreen] (z3) -- (ok);
\end{tikzpicture}
\end{center}
\vfill
\begin{itemize}
  \item Z3 is powerful but not infallible (solver bugs, bit-blasting limits)
  \item Excellent for arithmetic, array bounds, linear constraints
  \item \jugmark{One step below a full kernel proof}
\end{itemize}
\end{frame}

% ---------------------------------------------------------------------------
% Slide 48: Trust Level — RUNTIME_WITNESSED
% ---------------------------------------------------------------------------
\begin{frame}
\frametitle{Trust Level --- \texttt{RUNTIME\_WITNESSED}}
\begin{block}{Tests Ran and Passed}
Dynamic evidence: the claim held for \emph{all tested inputs}.\\
Good evidence, but \textbf{only covers the tested cases}.
\end{block}
\vspace{0.3cm}
\begin{center}
\begin{tikzpicture}[
  box/.style={draw, thick, rounded corners, minimum width=2.6cm,
              minimum height=0.8cm, font=\small, align=center}
]
  \node[box, fill=jugeoLight]                          (claim) {Claim};
  \node[box, fill=jugeoOrange!15, right=1.3cm of claim]  (run)  {Run 1000\\test cases};
  \node[box, fill=jugeoBlue!15, right=1.3cm of run]      (ok)   {\textcolor{jugeoBlue}{\checkmark}~\texttt{WITNESSED}};

  \draw[-{Stealth}, thick, jugeoBlue] (claim) -- (run);
  \draw[-{Stealth}, thick, jugeoBlue] (run) -- (ok);
\end{tikzpicture}
\end{center}
\vfill
\begin{itemize}
  \item ``I checked 1\,000 cases and they all worked''
  \item Does \textbf{not} guarantee the 1\,001st case
  \item Useful when solvers time out or formal proof is too expensive
  \item \jugmark{Honest about its limits --- unlike ``all tests pass = correct''}
\end{itemize}
\end{frame}

% ---------------------------------------------------------------------------
% Slide 49: Trust Level — COPILOT_SUGGESTED
% ---------------------------------------------------------------------------
\begin{frame}
\frametitle{Trust Level --- \texttt{COPILOT\_SUGGESTED}}
\begin{alertblock}{An AI Suggested This Is Correct}
A language model (e.g.\ GitHub Copilot) proposed the claim or invariant.\\
Useful hint, but \textbf{NOT} a proof.
\end{alertblock}
\vspace{0.3cm}
\begin{center}
\begin{tikzpicture}[
  box/.style={draw, thick, rounded corners, minimum width=2.8cm,
              minimum height=0.8cm, font=\small, align=center}
]
  \node[box, fill=jugeoOrange!15]                        (ai) {Copilot suggests\\invariant};
  \node[box, fill=jugeoOrange!10, right=1.5cm of ai]     (tag) {\texttt{COPILOT\_}\\
                                                                  \texttt{SUGGESTED}};
  \node[right=0.8cm of tag, font=\small, text=jugeoRed, align=left]
    {Hard trust\\ceiling};

  \draw[-{Stealth}, thick, jugeoOrange] (ai) -- (tag);
\end{tikzpicture}
\end{center}
\vfill
\begin{itemize}
  \item AI contributions are \jugmark{tracked, not banned}
  \item Cannot be silently promoted to \texttt{VERIFIED\_PROOF}
  \item Must be independently confirmed to raise trust level
\end{itemize}
\end{frame}

% ---------------------------------------------------------------------------
% Slide 50: Trust Level — The Lower Levels
% ---------------------------------------------------------------------------
\begin{frame}
\frametitle{Trust Levels --- The Lower Tiers}
\begin{columns}[T]
\begin{column}{0.32\textwidth}
\begin{block}{\texttt{ORACLE\_PROPOSED}}
  \begin{itemize}
    \item A human expert asserted this
    \item Reviewed but not mechanically checked
    \item ``Trust me, I'm a domain expert''
  \end{itemize}
\end{block}
\end{column}
\begin{column}{0.32\textwidth}
\begin{block}{\texttt{UNVERIFIED}}
  \begin{itemize}
    \item No evidence yet
    \item Default state for any new claim
    \item ``We haven't looked at this''
  \end{itemize}
\end{block}
\end{column}
\begin{column}{0.32\textwidth}
\begin{alertblock}{\texttt{CONTRADICTED}}
  \begin{itemize}
    \item Evidence \textbf{against} the claim
    \item A test failed, a counterexample found
    \item \textcolor{jugeoRed}{Cannot be ignored}
  \end{itemize}
\end{alertblock}
\end{column}
\end{columns}
\vfill
\begin{center}
\small \texttt{CONTRADICTED} is special: it \textbf{poisons} any composition it enters.\\
\jugmark{You cannot hide a known failure.}
\end{center}
\end{frame}

% ---------------------------------------------------------------------------
% Slide 51: The Trust Algebra — Formal Definition
% ---------------------------------------------------------------------------
\begin{frame}
\frametitle{The Trust Algebra --- Formal Definition}
\begin{center}
\vspace{0.2cm}
{\Large $\mathcal{T} = (E_{\mathrm{adm}},\;\leq,\;\oplus,\;\ominus,\;\uparrow,\;\downarrow)$}
\end{center}
\vspace{0.4cm}
\begin{center}
\small\textit{Don't panic!  We'll explain each piece.}
\end{center}
\vspace{0.3cm}
\begin{center}
\begin{tabular}{@{} c l @{}}
  \toprule
  \textbf{Symbol} & \textbf{Meaning} \\
  \midrule
  $E_{\mathrm{adm}}$ & The seven admissible trust levels \\
  $\leq$              & Ordering: \texttt{CONTRADICTED} $<$ \ldots\ $<$ \texttt{VERIFIED\_PROOF} \\
  $\oplus$            & \jugmark{Compose} evidence (take the weaker) \\
  $\ominus$           & \jugmark{Challenge} evidence (may demote) \\
  $\uparrow$          & \jugmark{Promote} (explicit, audited, never silent) \\
  $\downarrow$        & \jugmark{Demote} (when evidence is weakened) \\
  \bottomrule
\end{tabular}
\end{center}
\vfill
\begin{center}
\small Six operators, three laws, one guarantee: \jugmark{honest accounting}.
\end{center}
\end{frame}

% ---------------------------------------------------------------------------
% Slide 52: Operator — Composition (⊕)
% ---------------------------------------------------------------------------
\begin{frame}
\frametitle{Operator --- Composition ($\oplus$)}
\begin{block}{Rule: Take the Weaker Level}
When combining evidence from two channels, the result is the \textbf{minimum}.
\end{block}
\vspace{0.3cm}
\begin{center}
\begin{tikzpicture}[
  box/.style={draw, thick, rounded corners, minimum width=3.2cm,
              minimum height=0.7cm, font=\small, align=center}
]
  \node[box, fill=jugeoGreen!15]                       (a) {\texttt{SOLVER\_DISCHARGED}};
  \node[box, fill=jugeoOrange!15, below=0.5cm of a]   (b) {\texttt{COPILOT\_SUGGESTED}};
  \node[font=\Large, right=1cm of $(a.east)!0.5!(b.east)$] (op) {$\oplus$};
  \node[box, fill=jugeoOrange!15, right=0.6cm of op]  (r) {\texttt{COPILOT\_SUGGESTED}};

  \draw[-{Stealth}, thick, jugeoBlue] (a.east) -| ([xshift=-4pt]op.west);
  \draw[-{Stealth}, thick, jugeoBlue] (b.east) -| ([xshift=-4pt]op.west);
  \draw[-{Stealth}, thick, jugeoOrange] (op) -- (r);
\end{tikzpicture}
\end{center}
\vfill
\begin{center}
\Large\itshape ``A chain is only as strong as its weakest link.''
\end{center}
\end{frame}

% ---------------------------------------------------------------------------
% Slide 53: Why Weakest-Link?
% ---------------------------------------------------------------------------
\begin{frame}
\frametitle{Why Weakest-Link?}
\begin{exampleblock}{Scenario}
Z3 proved the \textbf{arithmetic} is correct.\\
A Copilot suggested the \textbf{loop invariant}.\\[4pt]
Your overall confidence is limited by the Copilot suggestion.
\end{exampleblock}
\vspace{0.4cm}
\begin{center}
\begin{tikzpicture}[
  box/.style={draw, thick, rounded corners, minimum width=2.4cm,
              minimum height=0.7cm, font=\small, align=center}
]
  \node[box, fill=jugeoGreen!15]                        (z3)  {Z3: arithmetic};
  \node[box, fill=jugeoOrange!15, right=1cm of z3]      (cop) {AI: invariant};
  \node[box, fill=jugeoOrange!15, right=1cm of cop]     (all) {Combined};

  \draw[-{Stealth}, thick, jugeoBlue]   (z3)  -- (cop);
  \draw[-{Stealth}, thick, jugeoOrange] (cop) -- (all);
\end{tikzpicture}
\end{center}
\vfill
\begin{itemize}
  \item \textcolor{jugeoRed}{\textbf{You cannot launder low-trust evidence through high-trust channels}}
  \item The weakest link determines overall strength
  \item This is by design --- it keeps the system \jugmark{honest}
\end{itemize}
\end{frame}

% ---------------------------------------------------------------------------
% Slide 54: Law 1 — No Silent Promotion
% ---------------------------------------------------------------------------
\begin{frame}
\frametitle{Law 1 --- No Silent Promotion}
\begin{alertblock}{The Rule}
You \textbf{CANNOT} silently upgrade trust.  Every promotion requires:
\end{alertblock}
\vspace{0.3cm}
\begin{center}
\begin{tikzpicture}[
  req/.style={draw, thick, rounded corners, fill=jugeoLight,
               minimum width=4.5cm, minimum height=0.65cm, font=\small}
]
  \node[req]                        (j) {\textbf{1.}~Explicit justification};
  \node[req, below=0.3cm of j]     (p) {\textbf{2.}~Named policy route};
  \node[req, below=0.3cm of p]     (a) {\textbf{3.}~Entry in append-only audit log};

  \draw[-{Stealth}, thick, jugeoBlue] (j) -- (p);
  \draw[-{Stealth}, thick, jugeoBlue] (p) -- (a);
\end{tikzpicture}
\end{center}
\vfill
\begin{itemize}
  \item Attempting silent promotion $\Rightarrow$ \textcolor{jugeoRed}{\texttt{JuGeoError}}
  \item The audit log is \textbf{append-only}: no rewriting history
  \item \jugmark{Trust must be earned, not assigned}
\end{itemize}
\end{frame}

% ---------------------------------------------------------------------------
% Slide 55: Code Example — Silent Promotion Fails
% ---------------------------------------------------------------------------
\begin{frame}[fragile]
\frametitle{Silent Promotion --- Fails}
\begin{lstlisting}[escapechar=@]
from jugeo.trust import TrustProfile, COPILOT_SUGGESTED, VERIFIED_PROOF

profile = TrustProfile(
    tier=COPILOT_SUGGESTED,
    source="copilot-invariant-suggestion",
    claim="loop_invariant_holds",
)

# Attempt silent promotion
profile.promote(VERIFIED_PROOF)
#  @\textcolor{jugeoRed}{\textrightarrow\ raises JuGeoError:}@
#  "Silent promotion from COPILOT_SUGGESTED
#   to VERIFIED_PROOF is forbidden.
#   Provide explicit=True with a reason."
\end{lstlisting}
\vfill
\begin{center}
\small The system \textbf{refuses} to upgrade trust without justification.\\
\jugmark{No shortcuts.}
\end{center}
\end{frame}

% ---------------------------------------------------------------------------
% Slide 56: Code Example — Explicit Promotion Works
% ---------------------------------------------------------------------------
\begin{frame}[fragile]
\frametitle{Explicit Promotion --- Works}
\begin{lstlisting}[escapechar=@]
from jugeo.trust import TrustProfile
from jugeo.trust import COPILOT_SUGGESTED, SOLVER_DISCHARGED

profile = TrustProfile(
    tier=COPILOT_SUGGESTED,
    source="copilot-invariant-suggestion",
)

# Explicit promotion with full justification
promoted = profile.promote(
    SOLVER_DISCHARGED,
    explicit=True,
    reason="z3-confirmed-arithmetic",
)
# @\textcolor{jugeoGreen}{\checkmark}@ promoted.tier == SOLVER_DISCHARGED
# @\textcolor{jugeoGreen}{\checkmark}@ Audit trail records promotion + reason
\end{lstlisting}
\vfill
\begin{center}
\small Every promotion is \jugmark{justified, named, and logged}.
\end{center}
\end{frame}

% ---------------------------------------------------------------------------
% Slide 57: Law 2 — Conservative Join
% ---------------------------------------------------------------------------
\begin{frame}
\frametitle{Law 2 --- Conservative Join}
\begin{block}{Rule: Always Take the Weaker Trust Level}
$a \oplus b = \min(a, b)$. This prevents \textbf{trust laundering}.
\end{block}
\vspace{0.2cm}
\begin{center}
\small
\begin{tabular}{@{} l c l c l @{}}
  \toprule
  \textbf{Channel A} & $\oplus$ & \textbf{Channel B} & $=$ & \textbf{Result} \\
  \midrule
  \texttt{VERIFIED\_PROOF}    & $\oplus$ & \texttt{VERIFIED\_PROOF}    & $=$ & \texttt{VERIFIED\_PROOF} \\
  \texttt{SOLVER\_DISCHARGED} & $\oplus$ & \texttt{RUNTIME\_WITNESSED} & $=$ & \texttt{RUNTIME\_WITNESSED} \\
  \texttt{SOLVER\_DISCHARGED} & $\oplus$ & \texttt{COPILOT\_SUGGESTED} & $=$ & \texttt{COPILOT\_SUGGESTED} \\
  \texttt{VERIFIED\_PROOF}    & $\oplus$ & \texttt{CONTRADICTED}       & $=$ & \texttt{CONTRADICTED} \\
  \bottomrule
\end{tabular}
\end{center}
\vfill
\begin{itemize}
  \item \texttt{CONTRADICTED} poisons everything --- you cannot ignore failure
  \item Composition is commutative: $a \oplus b = b \oplus a$
  \item \jugmark{No amount of strong evidence cancels a known contradiction}
\end{itemize}
\end{frame}

% ---------------------------------------------------------------------------
% Slide 58: Law 3 — Challenge Conservativity
% ---------------------------------------------------------------------------
\begin{frame}
\frametitle{Law 3 --- Challenge Conservativity}
\begin{alertblock}{Rule: Challenges May Demote, Never Ignore}
When evidence is challenged (e.g.\ a test fails), the system \textbf{must} respond.\\
Old trust cannot stand without explanation.
\end{alertblock}
\vspace{0.3cm}
\begin{center}
\begin{tikzpicture}[
  box/.style={draw, thick, rounded corners, minimum width=2.5cm,
              minimum height=0.7cm, font=\small, align=center}
]
  \node[box, fill=jugeoGreen!15]                         (before) {\texttt{SOLVER\_}\\
                                                                    \texttt{DISCHARGED}};
  \node[box, fill=jugeoRed!15, right=1cm of before]      (ch)     {Test fails\\($\ominus$)};
  \node[box, fill=jugeoOrange!15, right=1cm of ch]       (after)  {\texttt{RUNTIME\_}\\
                                                                    \texttt{WITNESSED}};

  \draw[-{Stealth}, thick, jugeoBlue] (before) -- (ch);
  \draw[-{Stealth}, thick, jugeoRed]  (ch) -- (after);
\end{tikzpicture}
\end{center}
\vfill
\begin{itemize}
  \item Like recalling a witness's testimony in court
  \item The demotion and its reason are logged in the audit trail
  \item \jugmark{The system never pretends a challenge didn't happen}
\end{itemize}
\end{frame}

% ---------------------------------------------------------------------------
% Slide 59: Why This Matters — A Real Scenario
% ---------------------------------------------------------------------------
\begin{frame}
\frametitle{A Real Scenario}
\begin{exampleblock}{Verifying \texttt{binary\_search}}
\small
\begin{description}[\textbf{Cover points:}]
  \item[\textbf{Cover points:}] 10 total obligations
  \item[\textbf{Z3 proved:}] 8 of 10 \hfill\textcolor{jugeoGreen}{\texttt{SOLVER\_DISCHARGED}}
  \item[\textbf{Z3 timed out:}] 2 edge cases \hfill\textcolor{jugeoOrange}{\texttt{---}}
  \item[\textbf{Runtime tests:}] confirmed the 2 edge cases \hfill\textcolor{jugeoBlue}{\texttt{RUNTIME\_WITNESSED}}
\end{description}
\end{exampleblock}
\vspace{0.3cm}
\begin{block}{JuGeo's Honest Report}
\begin{itemize}
  \item Overall trust: \texttt{RUNTIME\_WITNESSED} \hfill (conservative join)
  \item Residual obligations: \textbf{none}
  \item Every channel's contribution: \textbf{recorded}
\end{itemize}
\end{block}
\vfill
\begin{center}
\small Not ``fully proven'' --- but an \jugmark{honest, auditable account} of what was established and how.
\end{center}
\end{frame}

% ---------------------------------------------------------------------------
% Slide 60: Comparison — Trust Models
% ---------------------------------------------------------------------------
\begin{frame}
\frametitle{Trust Models --- Comparison}
\begin{center}
\begin{tabular}{@{} l c c c c @{}}
  \toprule
  & \textbf{LEAN} & \textbf{Coq} & \textbf{F*} & \jugmark{JuGeo} \\
  \midrule
  Trust levels       & 2 (binary)    & 2 (binary)    & 2 (binary)    & \textcolor{jugeoGreen}{\textbf{7}} \\
  Evidence channels  & 1             & 1             & 2 (SMT)       & \textcolor{jugeoGreen}{\textbf{6}} \\
  Audit trail        & \textcolor{jugeoRed}{No}  & \textcolor{jugeoRed}{No}  & \textcolor{jugeoRed}{No}  & \textcolor{jugeoGreen}{\textbf{Yes}} \\
  AI accounting      & \textcolor{jugeoRed}{No}  & \textcolor{jugeoRed}{No}  & \textcolor{jugeoRed}{No}  & \textcolor{jugeoGreen}{\textbf{Yes}} \\
  Challenge protocol & \textcolor{jugeoRed}{No}  & \textcolor{jugeoRed}{No}  & \textcolor{jugeoRed}{No}  & \textcolor{jugeoGreen}{\textbf{Yes}} \\
  Promotion control  & N/A           & N/A           & N/A           & \textcolor{jugeoGreen}{\textbf{Yes}} \\
  \bottomrule
\end{tabular}
\end{center}
\vfill
\begin{center}
\large JuGeo is the only system that gives you an\\
\jugmark{honest, multi-channel account of verification evidence}.
\end{center}
\end{frame}
